% Narrative literature review template (LaTeX)
% This is a lightweight starting point for surveys/related-work sections.

\documentclass[11pt]{article}

\usepackage[T1]{fontenc}
\usepackage[utf8]{inputenc}
\usepackage{lmodern}
\usepackage{microtype}
\usepackage{amsmath,amssymb}
\usepackage{graphicx}
\usepackage{booktabs}
\usepackage{hyperref}
\usepackage[numbers]{natbib}

\title{Title of Literature Review}
\author{Author Name}
\date{\today}

\begin{document}
\maketitle

\begin{abstract}
State the scope of the review, the selection methodology (briefly), and the main themes and gaps (150--250 words).
\end{abstract}

\section{Introduction}
Define the problem area and why it matters. Clarify the scope and what is out-of-scope.

\section{Method (Source Selection)}
Describe where you searched (databases), the time window, and the inclusion/exclusion criteria. This is required for defensible reviews.

\section{Thematic Synthesis}
Organize work by themes rather than by paper-by-paper summaries.

\subsection{Theme A: Foundational Approaches}
Summarize what is known and where results agree or disagree \citep{smith2023,jones2022}.
Discuss seminal contributions and how they shaped the field.
\citet{author2020} introduced the foundational framework that subsequent work builds upon.

\subsection{Theme B: Methodological Advances}
Compare approaches under consistent dimensions (assumptions, data, metrics, settings).
Recent studies \citep{other2021,diff2022} have proposed alternative methodologies with
varying degrees of success.

\subsection{Theme C: Applications and Evaluation}
Discuss how methods have been applied in practice and what evaluation criteria
are commonly used \citep{recent2023}.

% --- Comparison Table ---
\begin{table}[ht]
  \centering
  \caption{Comparison of Major Approaches}
  \label{tab:comparison}
  \begin{tabular}{llccc}
    \toprule
    Study & Approach & Dataset & Metric & Result \\
    \midrule
    \citet{smith2023}   & Method A & Dataset 1 & F1    & 0.87 \\
    \citet{jones2022}   & Method B & Dataset 1 & F1    & 0.83 \\
    \citet{other2021}   & Method C & Dataset 2 & Acc.  & 0.91 \\
    \citet{diff2022}    & Method D & Dataset 2 & Acc.  & 0.89 \\
    \bottomrule
  \end{tabular}
\end{table}

% --- Evidence Summary Table ---
\begin{table}[ht]
  \centering
  \caption{Evidence Summary Across Reviewed Studies}
  \label{tab:evidence}
  \begin{tabular}{lp{3cm}p{3cm}p{3cm}}
    \toprule
    Theme & Key Finding & Strength & Gaps \\
    \midrule
    Theme A & Established theoretical basis & Strong (multiple replications) & Limited to controlled settings \\
    Theme B & Improved efficiency & Moderate (few large-scale studies) & Scalability untested \\
    Theme C & Practical applicability & Emerging (case studies) & Lacks longitudinal data \\
    \bottomrule
  \end{tabular}
\end{table}

% --- Conceptual Framework Figure ---
\begin{figure}[ht]
  \centering
  % \includegraphics[width=0.7\textwidth]{figures/conceptual-framework.pdf}
  \caption{Conceptual framework illustrating the relationships between the
  three identified themes. Replace this placeholder with your diagram.}
  \label{fig:framework}
\end{figure}

\section{Open Problems and Gaps}
List gaps backed by evidence (e.g., missing ablations, weak external validity, lack of reproducibility artifacts).

\subsection{Research Gaps}
\begin{itemize}
  \item \textbf{Gap 1:} Limited cross-domain evaluation --- most studies validate on a single benchmark.
  \item \textbf{Gap 2:} Reproducibility concerns --- fewer than 30\% of reviewed papers release code or data.
  \item \textbf{Gap 3:} Lack of longitudinal studies --- current evidence is predominantly cross-sectional.
\end{itemize}

\section{Future Directions}
\begin{enumerate}
  \item Conduct large-scale, multi-site evaluations to improve external validity.
  \item Develop standardized benchmarks and reporting protocols.
  \item Investigate the long-term effects through longitudinal study designs.
\end{enumerate}

\section{Conclusion}
Summarize themes and propose future directions. Restate the most critical gaps
and their implications for the field.

\bibliographystyle{plainnat}
\bibliography{references}

\end{document}
